\providecommand{\main}{..}
\documentclass[\main/hippo]{subfiles}
\begin{document}

\section{Derivations of HiPPO Projection Operators}
\label{sec:derivations}

We derive the memory updates associated with the translated Legendre (LegT)
and translated Laguerre (LagT) measures as presented in
\cref{subsec:high_order_projection}, along with the scaling Legendre (LegS)
measure (\cref{sec:theory_legs}).
To show the generality of the framework, we also derive memory updates with
Fourier basis (recovering the Fourier Recurrent Unit~\citep{zhang2018learning})
and with Chebyshev basis.

The majority of the work has already been accomplished by setting up the
projection framework, and the proof simply requires following the technical
outline laid out in \cref{sec:hippo-framework-details}.
In particular, the definition of the coefficients \eqref{eq:coefficient} and reconstruction \eqref{eq:reconstruction} does not change,
and we only consider how to calculate the coefficients dynamics \eqref{eq:coefficient-dynamics}.

For each case, we follow the general steps:
\begin{description}%
  \item[Measure and Basis]
    define the measure $\mu^{(t)}$ or weight $\omega(t, x)$ and basis functions $p_n(t, x)$,
  \item[Derivatives]
    compute the derivatives of the measure and basis functions,
  \item[Coefficient Dynamics]
    plug them into the coefficient dynamics
    (equation~\eqref{eq:coefficient-dynamics}) to derive the ODE that describes how
    to compute the coefficients $c(t)$,
  \item[Reconstruction]
    provide the complete formula to reconstruct an approximation to the function $f_{\le t}$,
    which is the optimal projection under this measure and basis.
\end{description}

The derivations in \cref{sec:derivation-legt,sec:derivation-lagt}
prove \cref{thm:legt-lagt}, and the derivations in \cref{sec:derivation-legs}
prove \cref{thm:legs}.
\cref{sec:derivation-fourier,sec:derivation-chebyshev} show additional results for Fourier-based bases.

Figure~\ref{fig:measures} illustrates the overall framework when we use Legendre
and Laguerre polynomials as the basis,
contrasting our main families of time-varying measures $\mu^{(t)}$.



\begin{figure}
  \centering
  \includegraphics[width=\linewidth]{figures/measures.pdf}
  \caption{
    \textbf{Illustration of HiPPO measures.}
    At time $t_0$, the history of a function $f(x)_{x \le t_0}$ is summarized by polynomial approximation with respect to the measure $\mu^{(t_0)}$ (blue), and similarly for time $t_1$ (purple).
    (Left) The Translated Legendre measure (\textbf{LegT}) assigns weight in the window $[t-\theta, t]$.
    For small $t$, $\mu^{(t)}$ is supported on a region $x < 0$ where $f$ is not defined.
    When $t$ is large, the measure is not supported near $0$, causing the projection of $f$ to forget the beginning of the function.
    (Middle) The Translated Laguerre (\textbf{LagT}) measure decays the past exponentially.
    It does not forget, but also assigns weight on $x < 0$.
    (Right) The Scaled Legendre measure (\textbf{LegS}) weights the entire history $[0, t]$ uniformly.
  }
  \label{fig:measures}
\end{figure}



\IfFileExists{./derivations/legt.tex}{
  \subfile{./derivations/legt}

  \subfile{./derivations/lagt}

  \subfile{./derivations/legs}

  \subfile{./derivations/fourier}

  \subfile{./derivations/chebyshev}
  }{
  \subfile{./src/derivations/legt}
  \subfile{./src/derivations/lagt}
  \subfile{./src/derivations/legs}
  \subfile{./src/derivations/fourier}
  \subfile{./src/derivations/chebyshev}
}







\end{document}

